\chapter{Correspondenz}
\quellseite{31}

\begin{center}
Birmingham, den 5. Februar 1858.
\end{center}

Geehrtester Herr Professor!

Die Empfehlung von Augustins Bekenntnissen für die luth. theol. Pfarrerbibliothek im Decemberheft der „Lehre und Wehre“ veranlaßt mich zu der Anfrage, ob Sie nicht nachträglich eine ganz vorzügliche Ausgabe dieser trefflichen Schrift nennen wollen; ich meine: Sancti Augustini Confessionum libri tredecim, herausgegeben und erläutert von Karl von Raumer. Stuttgart bei S. G. Liesching, 1856. Zu Grunde gelegt ist die Oxforder Edition; die sprachlich schwierigen Stellen, und deren gibt es ja bei Augustin auch für einen geübten Lateiner nicht wenige, sind gut erklärt, der Inhalt selbst aber ist durch Parallelstellen aus den andern Schriften Augustins erläutert. Erst seit ich diese Ausgabe gebrauche, ist mir das köstliche Buch vollständig klar geworden und ich wünsche recht vielen Amtsbrüdern den gleichen Segen.

\begin{unterschrift}
Mit herzlichem Gruß und Segenswunsch Ihr\\
Herm. Gilbert.
\end{unterschrift}

\begin{center}
(Eingesandt.)
\end{center}
